\bookchapter{Chapitre 5 : Faire en sorte que votre smartphone vous comprenne mieux — paramètres de base}

\chapterintro{Quelques minutes de configuration améliorent considérablement l'expérience. Ces paramètres sont comme « l'accordage » du smartphone, le rendant plus conforme à vos habitudes d'utilisation.}

\booksection{Son et vibration}

Définissez une sonnerie et un son de notification que vous aimez. Utilisez flexiblement le mode silencieux (réunion), le mode Ne pas déranger (sommeil), le monde devient instantanément calme.

\begin{itemize}
  \item \textbf{Sonnerie}: Choisissez-en une que vous entendez et qui n'est pas gênante.
  \item \textbf{Mode vibration}: Alerte par vibration en mode silencieux.
  \item \textbf{Mode Ne pas déranger}: Des exceptions peuvent être définies, par exemple autoriser les appels des contacts.
  \item \textbf{Séparation des volumes}: Volume multimédia, volume de sonnerie, volume d'alarme peuvent être réglés séparément.
\end{itemize}

\booksection{Affichage et luminosité}

\begin{itemize}
  \item \textbf{Luminosité automatique}: Laissez le smartphone ajuster automatiquement selon la lumière ambiante, sans souci.
  \item \textbf{Mode protection des yeux/Nuit}: Filtre la lumière bleue, plus confortable pour regarder le smartphone le soir.
  \item \textbf{Mode sombre}: Économise la batterie (particulièrement efficace sur les écrans OLED), et n'éblouit pas la nuit. Beaucoup d'applications prennent aussi en charge le mode sombre.
  \item \textbf{Taux de rafraîchissement de l'écran}: Certains smartphones Android sont réglables, plus il est élevé (comme 120 Hz), plus le défilement est fluide.
  \item \textbf{Verrouillage automatique}: Réglez sur 15 ou 30 secondes, à la fois sécurisé et économe en énergie.
  \item \textbf{Taille de police}: Vous trouvez les caractères petits ? Agrandissez-les.
\end{itemize}

\booksection{Fond d'écran, thème et bureau}

Changez pour un beau fond d'écran (statique ou animé), l'humeur change aussi. Les smartphones Android de marque ont souvent une boutique de thèmes, permettant de changer l'ensemble des icônes et du style.

\textbf{Widgets de bureau}: Météo, agenda, contrôle musical sur le bureau, accès en une étape. iOS et Android le prennent tous deux en charge, appuyez longuement sur un espace vide du bureau pour ajouter.

\booksection{Paramètres d'efficacité : faites du smartphone votre « cerveau externe »}

Cette partie regroupe des fonctionnalités « dont on ne peut plus se passer une fois configurées », chacune pouvant améliorer significativement votre efficacité quotidienne.

\booksubsection{Mode Concentration / Ne pas déranger : reconquérez l'attention volée par les notifications}

\begin{itemize}
  \item \textbf{Idée centrale}: Pas « utiliser moins le smartphone », mais « se concentrer quand il faut se concentrer, être joignable quand il faut l'être ».
  \item \textbf{Conseils d'utilisation}:\newline · Mode travail : autoriser uniquement les contacts « favoris » et les applications nécessaires ;\newline · Mode sommeil : activation automatique, autoriser uniquement les alarmes ;\newline · Sport / Lecture / Réunion : personnaliser selon l'activité en cours.
  \item \textbf{Utilisation avancée}: Utilisez les automatisations + déclencheurs de localisation / temps pour que les modes basculent automatiquement. Par exemple : les raccourcis iOS permettent de configurer « en arrivant au bureau, basculer automatiquement en mode travail ».
\end{itemize}

\booksubsection{Gestion des notifications : que les pastilles rouges ne vous harcèlent plus}

\begin{itemize}
  \item \textbf{Désactiver les notifications marketing}: 90 \% des notifications push ne vous sont en fait pas utiles.
  \item \textbf{Classer par importance}: Mettez en haut les communications interpersonnelles (appels, SMS, messagerie instantanée) ; les actualités, shopping et réseaux sociaux peuvent être réglés en « notification silencieuse ».
  \item \textbf{Notifications groupées}: Repliez les notifications du même type ensemble, évitant d'être submergé par des dizaines de messages au déblocage.
\end{itemize}

\booksubsection{Gestes rapides : gagnez une longueur d'avance}

\begin{itemize}
  \item \textbf{Capture d'écran}: Bouton d'alimentation + volume bas (iOS plein écran) / Bouton d'alimentation + volume bas (Android universel).
  \item \textbf{Capture d'écran longue / capture défilante}: Après une capture normale, cliquez sur l'option « capture longue ».
  \item \textbf{Scan en un clic}: Ajoutez l'icône « scanner » dans le centre de contrôle / réglages rapides.
  \item \textbf{Toucher le dos (iPhone)}: Dans [Paramètres] > [Accessibilité] > [Toucher] > [Toucher le dos], associez « capture d'écran » ou « verrouillage » et d'autres actions.
  \item \textbf{Déclenchement vocal}: Siri / Xiao Ai / Xiao Yi / Xiao Bu, un seul mot pour vérifier la météo, régler une alarme.
\end{itemize}

\booksubsection{Clavier : l'outil d'efficacité le plus sous-estimé}

\begin{itemize}
  \item \textbf{Choisir le bon clavier}:\newline · Sogou / Baidu : dictionnaire étendu, mise à jour rapide des termes à la mode ;\newline · iFlytek : meilleur taux de reconnaissance vocale, outil parfait pour prendre des notes en réunion ;\newline · Gboard (Google) : simple, synchronisation multiplateforme puissante, saisie par glissement ;\newline · Clavier iOS intégré : meilleure confidentialité, intégration profonde au système.
  \item \textbf{Options indispensables}: Correction automatique, suggestion de mots, synchronisation du cloud, saisie de chiffres et symboles par glissement vers le haut.
  \item \textbf{Conseil de confidentialité}: Utilisez avec prudence les fonctions de « suggestion de phrase complète » ou de « recommandation personnalisée », qui peuvent télécharger le contenu de vos saisies.
\end{itemize}

\booksubsection{Bureau et multitâche : fluidifiez votre flux de travail}

\begin{itemize}
  \item \textbf{Écran partagé / Image dans l'image}: Répondre aux messages en regardant une vidéo, prendre des notes en naviguant sur le web.
  \item \textbf{Fenêtre flottante / fenêtre suspendue}: WeChat / TikTok prennent en charge le « mode fenêtre flottante », repliez après avoir répondu.
  \item \textbf{Applications fréquentes en haut}: La barre Dock (en bas de l'écran) ne contient que les 4 applications les plus utilisées, réduisant le temps de recherche.
\end{itemize}

\begin{tcbinfo}{Liste de contrôle « premier jour de travail »}
Après avoir reçu un nouveau smartphone ou une mise à jour majeure, il est recommandé de passer 10 minutes sur ces paramètres :\newline (1) Mot de passe d'écran de verrouillage + biométrie (2) Groupement des notifications + désactivation des pushes marketing (3) Taille de police / taux de rafraîchissement / mode sombre\newline (4) Clavier / navigateur / messagerie par défaut (5) Contact d'urgence / carte médicale (6) Paramètres de sauvegarde

\end{tcbinfo}

\booksection{Accessibilité et conception universelle : la force de la technologie bienveillante}

Ne vous laissez pas effrayer par le nom « accessibilité », ces fonctionnalités sont utiles pour tout le monde. Elles ont été initialement conçues pour les personnes en situation de handicap, mais elles sont aussi très utiles pour les personnes valides — mains temporairement occupées à porter un bébé, écouteurs en panne dans le métro, presbytie avec des petits caractères... Elles peuvent toutes vous dépanner. Au fond, la conception universelle illustre l'idée que\textbf{la technologie est bienveillante}: permettre à chacun de profiter également de la vie numérique.

\booksubsection{Les « yeux » des utilisateurs malvoyants — guide complet du lecteur d'écran}

Le lecteur d'écran est l'outil principal pour les utilisateurs malvoyants. Il\textbf{lit à voix haute tout le contenu visible}à l'écran, les utilisateurs pilotent le smartphone par gestes et commandes vocales. Voici les méthodes d'utilisation de VoiceOver iOS et TalkBack Android.

\begin{tcbinfo}{iOS VoiceOver}
\textbf{Chemin d'activation}: [Paramètres] > [Accessibilité] > [VoiceOver], activez le commutateur. Vous pouvez aussi triple-cliquer sur le bouton latéral (ou le bouton principal) pour basculer rapidement.

\textbf{Gestes de base}:

\begin{enumerate}
  \item Glisser à un doigt : parcourir le contenu de l'écran élément par élément ;
  \item Balayer à deux doigts : faire défiler la page entière haut/bas/gauche/droite ;
  \item Double tap à un doigt : activer l'élément en focus (équivalent à « cliquer ») ;
  \item Balayer à trois doigts : changer de page ;
  \item Triple tap à un doigt : agrandir ou sélectionner (différent du triple tap normal).
\end{enumerate}

\textbf{Paramètres personnalisés}: [VoiceOver] > [Voix], vous pouvez régler la\textbf{vitesse de lecture}(environ 50 \% par défaut, réglable à 70-80 \% pour une écoute plus fluide), le\textbf{ton}, changer de pack vocal. iOS intègre « Voix améliorée » et « Contrôleur vocal », le premier réduit les bruits de respiration et les pauses, pour une écoute plus naturelle.

\end{tcbinfo}

\begin{tcbinfo}{Android TalkBack}
\textbf{Chemin d'activation}: [Paramètres] > [Accessibilité] > [TalkBack], activez le commutateur. Ou appuyez longuement sur les touches de volume haut et bas pendant 3 secondes pour activer rapidement (certaines marques le prennent en charge).

\textbf{Gestes de base}:

\begin{enumerate}
  \item Glisser à un doigt gauche/droite : changer l'élément en focus ;
  \item Double tap à un doigt : activer l'élément ;
  \item Glisser à deux doigts haut/bas/gauche/droite : faire défiler ;
  \item Simple tap à trois doigts : entrer en mode lecture (lecture continue) ;
  \item Simple tap à quatre doigts : pause/reprise de la lecture.
\end{enumerate}

\textbf{Paramètres personnalisés}: [TalkBack] > [Paramètres de synthèse vocale], vous pouvez ajuster la\textbf{vitesse}, le\textbf{ton}, changer de\textbf{moteur vocal}(il est recommandé d'installer le moteur Google de synthèse vocale, lecture plus naturelle).

\end{tcbinfo}

\booksubsection{Rendre les petits caractères grands, le flou net — réglages visuels}

\booksubsection{Gestes de zoom}

\begin{itemize}
  \item \textbf{iOS}: [Paramètres] > [Accessibilité] > [Zoom], après activation,\textbf{tap à trois doigts sur l'écran}pour agrandir. Zoom plein écran par défaut, option de zoom fenêtre (fenêtre d'agrandissement local).\textbf{Pincez à deux doigts pour ajuster le facteur de zoom (1,5x à 15x)}, double tap à trois doigts pour verrouiller le rapport de zoom actuel.
  \item \textbf{Android}: [Paramètres] > [Accessibilité] > [Agrandissement], tap à trois doigts pour agrandir, re-tap à trois doigts pour quitter ; ou directement\textbf{triple tap à trois doigts sur l'écran}pour activer le zoom temporaire (maximum 8x).
\end{itemize}

\booksubsection{Texte à contraste élevé / Texte en gras}

\begin{itemize}
  \item \textbf{iOS}: [Paramètres] > [Accessibilité] > [Affichage et taille du texte], activez les commutateurs « Texte en gras », « Contraste amélioré », « Réduire la transparence », etc. L'« Inversion intelligente » permet de garder le texte noir sur fond clair, évitant le texte blanc sur fond blanc.
  \item \textbf{Android}: [Paramètres] > [Affichage] > [Taille de police], réglable sur quatre niveaux : « Petit, Normal, Grand, Très grand ». Huawei/Xiaomi offre en plus les commutateurs « Gras » et « Texte à contraste élevé ».
\end{itemize}

\booksubsection{Filtres de couleur (adapté au daltonisme)}

\begin{itemize}
  \item \textbf{iOS}: [Paramètres] > [Accessibilité] > [Affichage et taille du texte] > [Filtres de couleur], propose\textbf{quatre modes de daltonisme rouge/vert/bleu/jaune}, ainsi que les options « Nuances de gris » et « Inverser les couleurs ». Il est recommandé de faire d'abord un test de vision des couleurs (intégré au système) avant de choisir le filtre correspondant.
  \item \textbf{Android}: [Paramètres] > [Accessibilité] > [Correction des couleurs], inclut\textbf{trois modes de daltonisme}+ le commutateur « Inverser les couleurs ».
\end{itemize}

\booksubsection{Description d'images par IA : que l'IA soit vos « yeux »}

Quand le lecteur d'écran rencontre une image, l'approche traditionnelle est de la sauter ou de lire seulement le mot « image ». Mais désormais l'IA peut\textbf{décrire le contenu de l'image en une phrase}, très utile pour reconnaître les captchas, lire des images, identifier des menus.

\begin{tcbinfo}{Prise en charge de la description d'images par IA sur chaque système}
\begin{itemize}
  \item \textbf{iOS 15+}: Dans [Paramètres] > [Accessibilité] > [VoiceOver] > [Description vocale], activez « Description d'images ». Le système appelle un\textbf{modèle IA local}pour générer automatiquement une description textuelle des images (nécessite iPhone 12 et supérieur).
  \item \textbf{Nouvelle version de l'application TalkBack}: Google TalkBack intègre dans les versions récentes la fonction de\textbf{questions-réponses sur les images}, vous pouvez demander « qu'y a-t-il sur cette image », l'IA répond en temps réel (ne dépend pas strictement de la version Android, il suffit de mettre à jour l'application TalkBack).
  \item \textbf{Applications tierces}: « Seeing AI » (iOS), « Lookout » (Android), « Xiao Yi Bang Bang » (Huawei) sont d'excellents assistants visuels IA, capables de reconnaître des objets, lire des documents, identifier des expressions faciales.
\end{itemize}

\end{tcbinfo}

\begin{tcbpractice}{Scénario d'utilisation}
Essayez de demander à propos d'une photo de rue « Combien de personnes y a-t-il sur cette image ? Que font-elles ? », ou ouvrez une application de livraison et laissez l'IA lire les noms de plats sur le menu —\textbf{c'est vraiment « la technologie rendant la vie plus égalitaire »}.

\end{tcbpractice}

\booksubsection{Les « oreilles » des utilisateurs malentendants — guide complet des sous-titres en temps réel}

Les sous-titres en temps réel (Live Captions) sont une fonctionnalité\textbf{de transcription vocale au niveau système}, particulièrement utile pour les utilisateurs malentendants dans les environnements bruyants, les bibliothèques et les contextes de langue étrangère.

\begin{tcbinfo}{Activation et optimisation des sous-titres en temps réel}
\textbf{iOS}: [Paramètres] > [Accessibilité] > [Sous-titres en temps réel], activez le commutateur. Le système détecte tout l'audio de l'appareil (appels, vidéos, podcasts, FaceTime, mémos vocaux) et affiche automatiquement les sous-titres en haut de l'écran.

\textbf{Android 12+}: Appuyez sur la\textbf{touche de volume}cliquez sur l'icône « Sous-titres en temps réel » dans la fenêtre popup ; ou dans [Paramètres] > [Accessibilité] > [Sous-titres en temps réel], activez.

\textbf{Conseils d'optimisation de la précision}:

\begin{enumerate}
  \item Utilisez autant que possible dans un\textbf{environnement calme}, loin de la climatisation, du bruit des machines ;
  \item La reconnaissance est meilleure quand le locuteur\textbf{articule clairement et parle à vitesse modérée};
  \item Les conversations à plusieurs réduisent la précision, vous pouvez utiliser le mode « n'afficher qu'une seule personne » ;
  \item Prend en charge le mandarin, le cantonais, l'anglais et d'autres langues principales,\textbf{la reconnaissance des langues minoritaires est moins bonne};
  \item En cas de musique de fond ou de voix parasites en cours de parole, la reconnaissance peut être saccadée.
\end{enumerate}

\end{tcbinfo}

\booksubsection{Double canal pour les appels : voix vers texte + texte vers voix}

Quand l'utilisateur\textbf{ne peut ni entendre ni parler}(comme une aphonie soudaine, période de récupération après une opération de la gorge, bégaiement sévère), le smartphone peut servir de « messager ».

\begin{tcbinfo}{Transcription en temps réel des appels (TTY / RTT)}
\textbf{iOS}: [Paramètres] > [Accessibilité] > [RTT/TTY], après activation, un clavier s'affiche à l'écran pendant l'appel, la voix de l'interlocuteur est transcrite en texte en temps réel,\textbf{et ce que vous tapez est lu à voix haute}pour l'interlocuteur.

\textbf{Android}: [Paramètres] > [Accessibilité] > [Transcription en temps réel] + [Notifications sonores]. La « transcription en temps réel » affiche la parole de l'interlocuteur en texte ; les « notifications sonores » convertissent votre saisie en parole quand la voix n'est pas pratique.

\textbf{Compatibilité}:

\begin{itemize}
  \item \textbf{RTT (Real-Time Text)}: Les deux parties doivent le prendre en charge, le texte s'affiche au fur et à mesure de la frappe, sans délai, la plupart des opérateurs le prennent en charge ;
  \item \textbf{TTY (téléscripteur)}: Nécessite un appareil TTY dédié, progressivement remplacé par RTT ;
  \item La\textbf{fonction voix vers texte de WeChat/QQ}est aussi une alternative : appuyez longuement sur un message vocal « Convertir en texte ».
\end{itemize}

\end{tcbinfo}

\booksubsection{Fonctionnalités « petites mais utiles » pour tout le monde}

\begin{itemize}
  \item \textbf{Toucher le dos}(iOS 14+) : [Paramètres] > [Accessibilité] > [Toucher] > [Toucher le dos], vous pouvez associer le\textbf{double tap/triple tap}pour déclencher capture d'écran, verrouillage, centre de contrôle, etc. Très pratique quand vous portez des colis avec les deux mains.
  \item \textbf{AssistiveTouch / Boule flottante}: iOS [Accessibilité] > [Toucher] > [AssistiveTouch], Android [Accessibilité] > [Navigation flottante]. Bouton virtuel remplaçant les gestes ou les boutons physiques, dépanne quand les boutons sont défectueux ou pour une utilisation à une main.
  \item \textbf{Contrôle par commutateur}: Contrôlez le smartphone via des commutateurs externes, clic unique, mouvements de tête,\textbf{conçu pour les utilisateurs avec des handicaps moteurs}. iOS [Accessibilité] > [Contrôle par commutateur], Android [Accessibilité] > [Contrôle par commutateur].
  \item \textbf{Reconnaissance sonore}(iOS) : [Accessibilité] > [Reconnaissance sonore], après activation, le smartphone « écoute » la sonnette, les pleurs de bébé, les alarmes incendie, les sirènes, etc.,\textbf{et vibre + affiche une notification à l'écran lorsqu'un son est détecté}, très précieux pour les utilisateurs malentendants et pendant le sommeil.
  \item \textbf{Rappel par flash LED}: iOS [Accessibilité] > [Audio/Visuel] > [Rappel par flash LED], Android [Paramètres] > [Notifications] > [Rappel par flash]. Le flash clignote lors des appels et notifications,\textbf{permet de percevoir sans son même dans un environnement bruyant}.
  \item \textbf{Audio mono}: [Paramètres] > [Accessibilité] > [Audio] > [Audio mono], fusionne les canaux gauche et droit. Les utilisateurs avec une audition unilatérale altérée ne perdent pas la moitié des informations sonores.
\end{itemize}

\begin{tcbinfo}{La philosophie d'accessibilité d'Apple}
Apple investit massivement dans l'accessibilité, mettant à jour et ajoutant de nouvelles fonctionnalités d'accessibilité chaque année. Le site d'Apple propose une page dédiée à l'accessibilité, détaillant l'utilisation de chaque fonctionnalité. Les fonctionnalités d'accessibilité d'iOS sont généralement plus complètes et plus fluides que celles d'Android, c'est d'ailleurs une raison importante pour laquelle de nombreuses personnes en situation de handicap choisissent Apple.

\end{tcbinfo}

\begin{tcbpractice}{Prenez 30 secondes pour essayer}
Trouvez un moment libre et faites les expériences suivantes :

\begin{enumerate}
  \item Activez\textbf{VoiceOver / TalkBack}, essayez de naviguer dans WeChat en glissant à trois doigts sans utiliser vos yeux ;
  \item Ouvrez les\textbf{sous-titres en temps réel}, parlez dans le smartphone et vérifiez la précision de la reconnaissance ;
  \item Utilisez le\textbf{tap à trois doigts}pour agrandir un texte à 200 \%, vivez le quotidien d'un utilisateur malvoyant ;
  \item Désactivez toutes les fonctionnalités précédentes, et réfléchissez à « si j'avais besoin de ces fonctionnalités en permanence, comment je voudrais qu'elles soient conçues ».
\end{enumerate}

\end{tcbpractice}
